% SGA 3, Expose III, Loop 1 English reconstruction.
% Scope: Section 1 through Proposition 1.2.4(iii), stopping immediately
% before Section 1.3. The final item occupies the top of source p. 28.
% Authority: Polo--Gille Exp3-14oct24.pdf, component-local pages 23--28
% (p. 28 partial; combined-reader pages 133--138, p. 138 partial; printed
% pages 123--128). Force-OCR is a drafting locator; English wording was
% compared with Jacob Reinhold's CC BY 4.0 Markdown lineage at commit
% e7a259f3f8608ad3edf9bf6eead3fd504dd2d23e. The PDF controls prose,
% mathematics, numbering, notes, and diagram images.

\setcounter{footnote}{34}

% Source page: Exp3 component-local p. 23; combined-reader p. 133; printed p. 123.
\section*{\SGAInlineRefTarget{sga3:iii:section:1}1. 扩张与上同调}
\addcontentsline{toc}{section}{1. 扩张与上同调}

\subsection*{\SGAInlineRefTarget{sga3:iii:subsection:1-1}1.1.}
设 \(\mathscr C\) 是对纤维积稳定的范畴.
\footnote{编者注: \(\mathscr C\) 对纤维积稳定的假设是新增的.
在应用中, 当 \(\mathscr C\) 为概形范畴
\((\mathbf{Sch})\), 或为函子
\((\mathbf{Sch})^{\circ}\to(\mathbf{Ens})\) 的范畴时, 这一假设成立.
若不作这一假设, 则必须在下面改为假设纤维积
\(G\times_S\cdots\times_S G\) 可表.}
设 \(S\) 是 \(\mathscr C\) 的对象, \(G\) 是可表的
\(S\)-群, \(F\) 是取值于交换群且带有 \(G\) 作用的 \(S\)-函子.
上同调群 \(H^n(G,F)\) 已在 \SGARefLink{sga3:i:subsection:5-1}{I~5.1} 中定义.
回忆它们是记为 \(C^\ast(G,F)\) 的一个复形的上同调群. 记
\[
 (G/S)^n=G\times_S\cdots\times_S G
 \qquad(n\text{ 个因子}),
\]
则有
\[
 C^n(G,F)=\operatorname{Hom}_S((G/S)^n,F).
\]
由于 \(G\) 可表, 因而每个 \((G/S)^n\) 也可表, 故还有
\[
 C^n(G,F)=F((G/S)^n).
\]
由此及余边界算子的定义可知, 复形 \(C^\ast(G,F)\) 仅依赖于
\(F\) 在 \(\mathscr C/S\) 的某个满子范畴上的限制,
其对象是 \(G\) 在 \(S\) 上的笛卡尔幂. 因此:

\medskip
\noindent\SGAInlineRefTarget{sga3:iii:lemma:1-1-1}\textbf{引理 1.1.1.}
\textit{设 \(\mathscr C\) 是对纤维积稳定的范畴,
\(S\) 是 \(\mathscr C\) 的对象, \(G\) 是可表的
\(S\)-群. 用 \(\mathscr C(G)\) 表示 \(\mathscr C/S\)
中以 \(G\) 在 \(S\) 上的笛卡尔幂为对象的满子范畴.
设 \(F\) 和 \(F'\) 是两个取值于交换群且带有 \(G\) 作用的
\(S\)-函子. 若 \(F\) 和 \(F'\) 在 \(\mathscr C(G)\)
上的限制相同, 则有典范同构}
\[
 H^\ast(G,F)\xrightarrow{\sim}H^\ast(G,F').
\]

\subsection*{\SGAInlineRefTarget{sga3:iii:subsection:1-1-2}1.1.2. 上同调与标量限制.\protect\footnotemark}
\footnotetext{编者注: 本段的标题是新增的.}
再记录一个比较结果. 设 \(T\to S\) 是 \(\mathscr C\) 中的态射.
若 \(F\) 是取值于交换群的 \(T\)-函子,
则由标量限制得到的函子 (参见 \SGARefLink{sga3:ii:section:1}{第 II 篇, 1}),
\[
 F_1=\prod_{T/S}F,
\]
是取值于交换群的 \(S\)-函子, 并且有取值于群的
\(S\)-函子之间的态射\footnotemark
\[
 u:\prod_{T/S}
 \underline{\operatorname{Aut}}_{T\text{-gr.}}(F)
 \longrightarrow
 \underline{\operatorname{Aut}}_{S\text{-gr.}}(F_1).
\]
\footnotetext{编者注: 明确地说, 设 \(S'\to S\) 及
\(\alpha\in
\underline{\operatorname{Aut}}_{(T\times_S S')\text{-gr.}}
(F_{T\times_S S'})\). 对每个 \(U\to T\times_S S'\),
设 \(\alpha_U\in
\operatorname{Aut}_{\mathrm{gr.}}(F_{T\times_S S'}(U))\)
是由 \(\alpha\) 定义的元素. 则 \(u(S')(\alpha)\) 是
\[
 \underline{\operatorname{Aut}}_{S\text{-gr.}}(F_1)(S')
 =
 \underline{\operatorname{Aut}}_{S'\text{-gr.}}
 \left(\prod_{T\times_S S'/S'}F_{T\times_S S'}\right)
\]
中对每个 \(S''\to S'\) 都满足
\(\beta_{S''}=\alpha_{T\times_S S''}\) 的元素 \(\beta\).}

现设 \(G\) 是取值于群的 \(S\)-函子, 并假设由
\[
 G_T\longrightarrow
 \underline{\operatorname{Aut}}_{T\text{-gr.}}(F).
\]
给定了 \(G_T\) 在 \(F\) 上的作用.
由 \(\prod_{T/S}\) 的定义, 这给出取值于群的 \(S\)-函子之间的态射
\[
 G\longrightarrow
 \prod_{T/S}\underline{\operatorname{Aut}}_{T\text{-gr.}}(F).
\]
因此, 与 \(u\) 复合就给出 \(G\) 在
\(F_1=\prod_{T/S}F\) 上的作用.
\footnote{编者注: 明确地说, 对每个 \(S'\to S\),
\(G(S')\) 在 \(F_1(S')=F(T\times_S S')\) 上的作用由
\(G_T(T\times_S S')=G(T\times_S S')\) 在 \(F(T\times_S S')\)
上的作用, 连同与投影 \(T\times_S S'\to S'\)
相应的群态射 \(G(S')\to G(T\times_S S')\) 给出.
等价地, \(G_T\) 在 \(F\) 上的作用是一个态射
\(G_T\times_T F\to F\). 应用 \(\prod_{T/S}\), 并记
\(G_1=\prod_{T/S}G_T\), 由 \SGARefLink{sga3:ii:remark:1-2}{II~1.2}
得到态射 \(G_1\times_S F_1\to F_1\), 因而得到
\(G_1\) 在 \(F_1\) 上的作用. \(G\) 的作用通过自然态射
\(G\to G_1\) 得到, 此态射通过伴随对应于
\(\operatorname{id}_G\).}

% Source page: Exp3 component-local p. 24; combined-reader p. 134; printed p. 124.
\medskip
\noindent\SGAInlineRefTarget{sga3:iii:lemma:1-1-2}\textbf{引理 1.1.2.}
\textit{在前述假设下, 有典范同构}
\[
 H^\ast\!\left(G,\prod_{T/S}F\right)
 \simeq H^\ast(G_T,F).
\]
事实上, 由上同调的定义, 标准复形典范同构.

\subsection*{\SGAInlineRefTarget{sga3:iii:subsection:1-2}1.2. 群态射的提升.\protect\footnotemark}
\footnotetext{编者注: 此标题是新增的.}
遵循一般原则, 作如下定义.

\medskip
\noindent\SGAInlineRefTarget{sga3:iii:definition:1-2-1}\textbf{定义 1.2.1.}
\textit{设}
\[
 1\longrightarrow M\xrightarrow{u}E\xrightarrow{v}G
\]
\textit{是 \(\widehat{\mathscr C}\)-群的态射列.
若下列等价条件成立, 则称此列\emph{正合 (exact)}:}
\begin{enumerate}
 \item[(i)] \textit{对每个 \(S\in\operatorname{Ob}\mathscr C\),
 下列通常意义下的群的序列正合:}
 \[
  1\longrightarrow M(S)\xrightarrow{u(S)}E(S)
   \xrightarrow{v(S)}G(S);
 \]

 \item[\SGAInlineRefTarget{sga3:iii:definition:1-2-1:item:ii}(ii)]
 \textit{对 \(\widehat{\mathscr C}\) 的每个对象 \(H\),
 下列通常意义下的群的序列正合:}
 \[
  1\longrightarrow\operatorname{Hom}(H,M)
   \xrightarrow{u(H)}\operatorname{Hom}(H,E)
   \xrightarrow{v(H)}\operatorname{Hom}(H,G).
 \]
\end{enumerate}

在 \SGARefLink{sga3:iii:definition:1-2-1:item:ii}{(ii)} 中取 \(H=G\),
可知 \(v\) 的截面 (未必保持群结构) 所成的集合要么为空,
要么在 \(\operatorname{Hom}(G,M)\) 下主齐性. 假定它非空, 并设
\[
 s:G\longrightarrow E
\]
是 \(v\) 的一个截面. 对每个 \(S\in\operatorname{Ob}\mathscr C\)
及每个 \(x\in G(S)\), 元素 \(s(x)\in E(S)\) 定义了
\(E_S\) 的一个内自同构, 它正规化 \(M_S\)
(更确切地说, 正规化 \(M_S\) 在 \(u_S\) 下的像),
因而定义了 \(M_S\) 的一个自同构.

% Source page: Exp3 component-local p. 25; combined-reader p. 135; printed p. 125.
\medskip
\noindent\SGAInlineRefTarget{sga3:iii:scholium:1-2-1-1}\textbf{附论 1.2.1.1.}
\footnote{编者注: 新增编号 \SGARefLink{sga3:iii:scholium:1-2-1-1}{1.2.1.1}
和 \SGARefLink{sga3:iii:definition:1-2-1-2}{1.2.1.2},
以突出此处引入的概念.}
\textit{若 \(M\) 交换, 从集合论角度可见, 这一自同构与所选截面无关,
仅依赖于 \(x\), 且关于 \(x\) 是乘法性的. 总之, 每个
\(M\) 交换且 \(v\) 容许截面的正合列}
\[
 (E)\qquad
 1\longrightarrow M\xrightarrow{u}E\xrightarrow{v}G
\]
\textit{都确定一个 \(\widehat{\mathscr C}\)-群态射}
\[
 G\longrightarrow
 \underline{\operatorname{Aut}}_{\widehat{\mathscr C}\text{-gr.}}(M),
\]
\textit{称为\emph{由扩张 \((E)\) 定义的 \(G\) 在 \(M\) 上的作用}.}

\medskip
\noindent\SGAInlineRefTarget{sga3:iii:definition:1-2-1-2}\textbf{定义 1.2.1.2.}
我们在 \SGARefLink{sga3:i:proposition:2-3-7}{I~2.3.7} 中看到,
\(v\) 容许作为 \(\widehat{\mathscr C}\)-群态射的截面,
当且仅当扩张 \((E)\) 作为扩张同构于 \(M\) 与 \(G\)
关于前述作用的半直积. 这样的 \(v\) 的截面称为
\emph{扩张 \((E)\) 的截面}, 简称 \emph{\((E)\) 的截面}.

若 \(s\) 是 \((E)\) 的截面, 且
\[
 m\in\Gamma(M)\simeq
 \operatorname{Ker}\bigl(\Gamma(E)\to\Gamma(G)\bigr)
\]
(\(\Gamma\) 的定义见 \SGARefLink{sga3:i:subsection:1-2}{I~1.2}), 则由
\[
 x\longmapsto u(m)s(x)u(m)^{-1}
\]
定义的态射 \(G\to E\) 也是 \((E)\) 的截面.
\footnote{编者注: 右端的 \(x\) 已改为 \(s(x)\).}
称它为由 \(m\), 或由 \(u(m)\), 所定义的内自同构从 \(s\) 导出的截面.

\medskip
\noindent\SGAInlineRefTarget{sga3:iii:lemma:1-2-2}\textbf{引理 1.2.2.}
\textit{设}
\[
 (E)\qquad
 1\longrightarrow M\xrightarrow{u}E\xrightarrow{v}G
\]
\textit{是 \(\widehat{\mathscr C}\)-群的正合列,
其中 \(M\) 交换, 且 \(v\) 容许截面.
令 \(G\) 按 \((E)\) 所定义的方式作用于 \(M\).}
\begin{enumerate}
 \item[\SGAInlineRefTarget{sga3:iii:lemma:1-2-2:item:i}(i)]
 \textit{扩张 \((E)\) 典范地定义一个类
 \(c(E)\in H^2(G,M)\), 此类为零是 \((E)\) 存在截面的充要条件.}

 \item[\SGAInlineRefTarget{sga3:iii:lemma:1-2-2:item:ii}(ii)]
 \textit{若 \(c(E)=0\), 则 \((E)\) 的截面集在 \(Z^1(G,M)\)
 下主齐性, 并且将 \((E)\) 的截面集模去由 \(\Gamma(M)\)
 的元素定义的内自同构的作用, 所得集合在 \(H^1(G,M)\) 下主齐性.}

 \item[\SGAInlineRefTarget{sga3:iii:lemma:1-2-2:item:iii}(iii)]
 \footnote{编者注: 此项是作为 \SGARefLink{sga3:iii:proposition:1-2-4:item:iii}{1.2.4(iii)}
 的对应项新增的. 此断言对 \(v\) 的每个截面都成立; 见证明.}
 \textit{设 \(s\) 是 \((E)\) 的截面. \(s\) 在由 \(\Gamma(M)\)
 定义的内自同构下的共轭所成的集合与 \(\Gamma(M)/H^0(G,M)\) 双射.}
\end{enumerate}

证明与通常群的情形完全相同; 从 \(v\) 的一个截面出发,
就保证了集合论计算的函子性. 我们简要指出其主要步骤.

% Source page: Exp3 component-local p. 26; combined-reader p. 136; printed p. 126.
\smallskip
\noindent\textbf{(a)}
对 \(v\) 的每个截面 \(s\), 配给一个态射
\[
 D_s:G\times G\longrightarrow M
\]
它从集合论角度由
\[
 u(D_s(x,y))=s(xy)s(y)^{-1}s(x)^{-1}.
\]
定义. 下面的计算表明 \(D_s\) 是一个
2-余循环.\footnote{编者注: 为与 \SGARefLink{sga3:i:subsection:5-1}{I~5.1}
相容, 已改正原文.} 由标准复形中微分的定义
(\SGARefLink{sga3:i:subsection:5-1}{I~5.1}),
\[
\begin{aligned}
 (\partial^2D_s)(x,y,z)
  ={}&\bigl(s(x)D_s(y,z)s(x)^{-1}\bigr)
     \cdot D_s(x,y)^{-1}\cdot D_s(xy,z)^{-1}\cdot D_s(x,yz).
\end{aligned}
\]
将 \(D_s\) 的定义代入此式, 不使用交换性就得到
\[
 D_s\in Z^2(G,M).
\]

\smallskip
\noindent\textbf{(b)}
若 \(s\) 和 \(s'\) 是 \(v\) 的两个截面, 则有态射
\(f:G\to M\), 使 \(s(x)=f(x)s'(x)\). 则
\[
\begin{aligned}
 D_{s'}(x,y)
 &=f^{-1}(xy)D_s(x,y)s(x)f(y)s(x)^{-1}f(x)\\
 &=D_s(x,y)\cdot f^{-1}(xy)
   \cdot\bigl(s(x)f(y)s(x)^{-1}\bigr)\cdot f(x),
\end{aligned}
\]
因而
\[
 D_{s'}=D_s\cdot\partial^1f.
\]
这说明 \(D_s\) 在 \(H^2(G,M)\) 中的类与所选的 \(v\)
的截面 \(s\) 无关; 它就是扩张 \((E)\) 的类 \(c(E)\).
\footnote{编者注: 前面这句话是新增的.}

\smallskip
\noindent\textbf{\SGAInlineRefTarget{sga3:iii:lemma:1-2-2:item:c}(c)}
设 \(s\) 和 \(s'\) 是 \(v\) 的两个截面, 并设
\(m\in\Gamma(M)\). 等式
\[
 s(x)=m^{-1}s'(x)m
 \qquad(x\in G(S),\ S\in\operatorname{Ob}\mathscr C)
\]
等价于
\[
 s(x)=m^{-1}s'(x)m\,s'(x)^{-1}s'(x),
 \qquad\text{即}\qquad
 s=\partial^0m\cdot s'.
\]
特别地, \(s\) 在 \(\Gamma(M)\) 中的稳定子是满足
\(\partial^0m=e_M\) 的元素 \(m\) 所成的子群, 即 \(H^0(G,M)\).
这就已经证明了 \SGARefLink{sga3:iii:lemma:1-2-2:item:iii}{(iii)}.

\smallskip
\noindent\textbf{(d)}
现在的论证是熟知的.\footnote{编者注: 下面扩充了原文.}
设 \(s_0\) 是 \(v\) 的任意截面. 存在一个必然形如
\(s=f\cdot s_0\) 的截面 \(s\) 为群态射, 即满足 \(D_s=0\),
当且仅当对某个 \(f:G\to M\) 有
\[
 (D_{s_0})^{-1}=\partial^1f,
\]
等价地, 当且仅当 \(c(E)=0\).
这就证明了 \SGARefLink{sga3:iii:lemma:1-2-2:item:i}{(i)}.

在此情形下, \((E)\) 的截面恰为截面
\(s'=h\cdot s\), 其中 \(h:G\to M\) 满足
\(\partial^1h=0\), 即 \(h\in Z^1(G,M)\).
此外, 由 \SGARefLink{sga3:iii:lemma:1-2-2:item:c}{(c)},
两个这样的截面 \(h_1\cdot s\) 和 \(h_2\cdot s\)
在 \(\Gamma(M)\) 下共轭, 当且仅当 \(h_1\) 和 \(h_2\)
在 \(H^1(G,M)\) 中的像相同. 这就证明了
\SGARefLink{sga3:iii:lemma:1-2-2:item:ii}{(ii)}.

仍设
\[
 (E)\qquad
 1\longrightarrow M\xrightarrow{u}E\xrightarrow{v}G
\]
是 \(M\) 交换的 \(\widehat{\mathscr C}\)-群正合列, 并设
\[
 f:H\longrightarrow G
\]
是 \(\widehat{\mathscr C}\)-群态射.

% Source page: Exp3 component-local p. 27; combined-reader p. 137; printed p. 127.
考虑 \(E_f=H\times_G E\). 它是
\(\widehat{\mathscr C}\)-群, 投影
\(v_f:E_f\to H\) 是 \(\widehat{\mathscr C}\)-群态射;
\(e_f:E_f\to E\) 也如此. 将 \(M\) 通过 \(u\) 映到 \(E\),
通过单位态射映到 \(H\), 也定义一个
\(\widehat{\mathscr C}\)-群态射 \(u_f:M\to E_f\).
因此我们构造了如下 \(\widehat{\mathscr C}\)-群的交换图.

\SGAThreeNativeDiagram[0.82\textwidth]
  {figures/exp3/diagram_028_exp3_p027_pullback_extension_diagram.tex}
  {拉回扩张 \((E_f)\) 及其到 \((E)\) 的态射.}
  {fig:exp3-pullback-extension}
  {Polo--Gille Exp3-14oct24.pdf, 分篇第 27 页; 合订本第 137 页; 印刷页码 127}

立即得到:

\medskip
\noindent\SGAInlineRefTarget{sga3:iii:lemma:1-2-3}\textbf{引理 1.2.3.}
\begin{enumerate}
 \item[(i)] \textit{列 \((E_f)\) 正合.}

 \item[(ii)] \textit{映射 \(s\mapsto e_f\circ s=f'\)
 给出满足 \(v_f\circ s=\operatorname{id}\) 的态射}
 \[
  s:H\longrightarrow E_f
 \]
 \textit{即 \(v_f\) 的截面与满足 \(v\circ f'=f\) 的态射}
 \[
  f':H\longrightarrow E
 \]
 \textit{即提升 \(f\) 的态射 \(f'\) 之间的一一对应.}

 \item[(iii)] \textit{在此对应下, \((E_f)\) 的截面
 对应于提升 \(f\) 的群态射 \(f'\).}
\end{enumerate}

将 \SGARefLink{sga3:iii:lemma:1-2-2}{引理~1.2.2} 应用于扩张 \((E_f)\),
并考虑 \SGARefLink{sga3:iii:lemma:1-2-3}{1.2.3}, 得到下面的命题,
它在形式上包含了 \SGARefLink{sga3:iii:lemma:1-2-2}{引理~1.2.2}.

\medskip
\noindent\SGAInlineRefTarget{sga3:iii:proposition:1-2-4}\textbf{命题 1.2.4.}
\textit{设}
\[
 (E)\qquad 1\longrightarrow M\longrightarrow E\xrightarrow{v}G
\]
\textit{是 \(M\) 交换的 \(\widehat{\mathscr C}\)-群正合列.
设 \(f:H\to G\) 是 \(\widehat{\mathscr C}\)-群态射,
并假设它提升为一个未必是群态射的态射 \(f':H\to E\).
令 \(H\) 通过复合态射}
\[
 H\xrightarrow{f'}E\xrightarrow{\operatorname{int}}
 \underline{\operatorname{Aut}}_{\widehat{\mathscr C}\text{-gr.}}(M).
\]
\textit{作用于 \(M\); 此复合态射是乘法性的, 且与 \(f'\) 的选择无关.}
\begin{enumerate}
 \item[(i)] \textit{态射 \(f\) 典范地定义一个类
 \(c(f)\in H^2(H,M)\), 此类为零是存在提升 \(f\)
 的 \(\widehat{\mathscr C}\)-群态射}
 \[
  f':H\longrightarrow E
 \]
 \textit{的充要条件.}

 \item[(ii)] \textit{若 \(c(f)=0\), 则提升 \(f\) 的
 \(\widehat{\mathscr C}\)-群态射 \(f'\) 所成的集合,
 模去由 \(\Gamma(M)\) 的元素 (即满足 \(v(m)=e\) 的元素
 \(m\in\Gamma(E)\)) 定义的内自同构的作用后,
 所得集合在 \(H^1(H,M)\) 下主齐性.}

 % Source page: Exp3 component-local p. 28 (top only); combined-reader
 % p. 138 (top only); printed p. 128.
 \item[\SGAInlineRefTarget{sga3:iii:proposition:1-2-4:item:iii}(iii)]
 \textit{若 \(f':H\to E\) 是提升 \(f\) 的群态射,
 则 \(f'\) 在由 \(\Gamma(M)\) 的元素定义的内自同构下的变换所成的集合同构于}
 \[
  \Gamma(M)/\Gamma(M^H)=\Gamma(M)/H^0(H,M).
 \]
\end{enumerate}

% Continue with Section 1.3 on the same source page.

% Translation coverage: entire component, Section 1 through Proposition 1.2.4(iii).
% Source lines: 441; translated body lines before this audit comment: 417.
% Math preserved: 239 inline formulas, 36 display blocks; two prose \text arguments translated.
% Reference targets: 18/18; reference links: 18/18; footnotes: 8/8 plus 3/3 footnotetext; diagrams: 1/1.
% Unresolved TODOs: 0; compile and visual review delegated to the volume integrator.
